\SourceLeaf{471}
\par\Needspace{12\baselineskip}\bigskip
\pdfbookmark[2]{Exposition, Quatrième Section, Erster Entwurf}{pallas-II-section4-first}
\begin{center}
{\large\bfseries Quatrième Section.\quad [Erster Entwurf.]\par}
\medskip
\textit{Principes de la détermination du mouvement pour un tems illimité.}
\end{center}
\pdfbookmark[3]{Exposition, [21]}{pallas-II-article21}
\begin{center}[21.]\end{center}
\noindent\hspace*{\parindent}L’intégration des variations instantanées des élémens fondée sur leur valeurs numériques ne donne les élémens que pour un tems limité. Pour pouvoir exécuter l’intégration indéfinie il faut auparavant exprimer ces variations en fonctions du tems telles qui soient susceptibles d’intégration; c’est en effet ce qui fait proprement la difficulté du problème. Vu la petitesse des masses perturbatrices comparées à la masse du soleil, nous considérons les variations des élémens, ou généralement les quantités qui contiennent le facteur \(\mu\) comme des quantités très petites du \textit{premier ordre}; celles qui en contiennent le quarré, ou le produit de \(\mu\) en une autre masse perturbatrice seront du second ordre et ainsi de suite. Désormais nous désignerons la vraie valeur de la longitude du noeud (comptée sur l’écliptique fixe) pour un instant indéterminé par \(\AscendingNode+\delta\AscendingNode\), de sorte que \(\AscendingNode\) soit une quantité constante et \(\delta\AscendingNode\) la partie variable du premier ordre; nous en userons de même pour les autres élémens. De là il est évident, que si l’on détermine la variation instantanée d’un élément en calculant toutes les quantités qui entrent dans sa valeur d’après les parties constantes, des élémens, cette variation sera approchée aux quantités du second ordre (exclus.); et généralement, que l’erreur des variations des élémens sera d’un ordre plus élevé d’un degré, que ne l’étaient les valeurs des élémens, qu’on avait employées dans son calcul. De là naît donc la méthode des approximations successives. Mais il faut avouer que toute simple que soit cette méthode dans l’idée, l’exécution en pourrait devenir presque impraticable, s’il fallait seulement calculer \textit{complétement} la seconde approximation. Heureusement pour les besoins de l’astronomie cela n’est nullement nécessaire; on se bornera à la première approximation, en y ajoutant quelques modifications tirées de la considération raisonnée des termes du second ordre.
